\documentclass[11pt, a4paper]{article}
\usepackage{amsmath, amssymb, bm}
\usepackage{geometry}
\usepackage{booktabs}
\usepackage{parskip}
\usepackage[hidelinks]{hyperref}
\usepackage{xcolor}
\usepackage{titlesec}

\geometry{margin=2.5cm}

\definecolor{accent}{RGB}{30, 140, 100}
\titleformat{\section}{\large\bfseries\color{accent}}{}{0em}{}[\titlerule]
\titleformat{\subsection}{\normalsize\bfseries}{}{0em}{}

\title{\textbf{Limb Align — Alignment Mathematics}\\[0.4em]
  \large Reference document for the auto-orientation pipeline}
\author{Limb Align}
\date{\today}

\begin{document}
\maketitle
\tableofcontents
\vspace{1em}

This document describes every mathematical operation performed when a user
clicks \textit{Auto Orient}.  The pipeline runs entirely in the browser using
JavaScript and Three.js.  Notation: bold lowercase $\bm{v}$ denotes a 3-D
column vector; bold uppercase $\bm{M}$ denotes a matrix.

% ─────────────────────────────────────────────────────────────────────────────
\section{Principal Component Analysis (PCA)}

\subsection{Centroid}

Given $n$ vertex positions $\bm{p}_1, \ldots, \bm{p}_n \in \mathbb{R}^3$, the
centroid is

\begin{equation}
  \bar{\bm{p}} = \frac{1}{n} \sum_{i=1}^{n} \bm{p}_i
\end{equation}

\subsection{Covariance matrix}

The $3 \times 3$ covariance matrix of the mean-subtracted vertex cloud is

\begin{equation}
  \bm{C} = \frac{1}{n} \sum_{i=1}^{n}
    (\bm{p}_i - \bar{\bm{p}})\,(\bm{p}_i - \bar{\bm{p}})^{\!\top}
  =
  \begin{pmatrix}
    c_{xx} & c_{xy} & c_{xz} \\
    c_{xy} & c_{yy} & c_{yz} \\
    c_{xz} & c_{yz} & c_{zz}
  \end{pmatrix}
\end{equation}

Because the matrix is real and symmetric, its eigenvectors are orthogonal and
its eigenvalues are real and non-negative.

\subsection{Power iteration — dominant eigenvector}

The dominant eigenvector (the direction of greatest variance, i.e.\ the
\emph{long axis} of the limb) is found by repeated multiplication and
re-normalisation, starting from $\bm{v}_0 = (0.577,\, 0.577,\, 0.577)^\top$:

\begin{equation}
  \bm{v}_{k+1} = \frac{\bm{C}\,\bm{v}_k}{\|\bm{C}\,\bm{v}_k\|},
  \qquad k = 0, 1, \ldots, 63
\end{equation}

After 64 iterations the sequence has converged to the dominant eigenvector
$\hat{\bm{l}}$ (unit long axis).

% ─────────────────────────────────────────────────────────────────────────────
\section{Proximal/Distal Detection}

\subsection{Projection onto the long axis}

Every vertex is projected onto the long axis to find the two extremal points:

\begin{equation}
  s_i = (\bm{p}_i - \bar{\bm{p}}) \cdot \hat{\bm{l}}
\end{equation}

\begin{equation}
  \bm{e}_A = \bm{p}_{\,\arg\min_i\, s_i},
  \qquad
  \bm{e}_B = \bm{p}_{\,\arg\max_i\, s_i}
\end{equation}

\subsection{Local radius at each endpoint}

The local radius at an endpoint $\bm{e}$ is the mean 3-D distance to all
vertices whose projection falls within 10\% of the total limb length of the
endpoint's projection:

\begin{equation}
  r(\bm{e}) = \frac{1}{|N(\bm{e})|}
              \sum_{i \in N(\bm{e})} \|\bm{p}_i - \bm{e}\|,
  \qquad
  N(\bm{e}) = \bigl\{\, i \;\big|\;
    \bigl| s_i - (\bm{e} \cdot \hat{\bm{l}}) \bigr| < 0.1\,(s_{\max} - s_{\min})
  \bigr\}
\end{equation}

\subsection{Distal end identification}

The distal end is the terminus with the \emph{smaller} local radius (anatomically,
the narrower end of a residual limb):

\begin{equation}
  \bm{e}_{\mathrm{distal}} =
  \begin{cases}
    \bm{e}_A & \text{if } r(\bm{e}_A) < r(\bm{e}_B) \\
    \bm{e}_B & \text{otherwise}
  \end{cases}
\end{equation}

\subsection{Orientation confidence}

The confidence metric quantifies how asymmetric the two end radii are.  A
limb that tapers strongly gives a high confidence score:

\begin{equation}
  c = \left(1 - \frac{\min(r_A,\, r_B)}{\max(r_A,\, r_B)}\right) \times 100\%
\end{equation}

% ─────────────────────────────────────────────────────────────────────────────
\section{Anatomical Axis Construction}

\subsection{Z-axis (proximal direction)}

The proximal $+Z$ axis points from the distal end toward the proximal end along
the long axis.  Starting from the raw long axis $\hat{\bm{l}}$:

\begin{equation}
  \hat{\bm{z}}_{\mathrm{raw}} =
  \begin{cases}
    \hat{\bm{l}}  & \text{if } \bm{e}_A \text{ is distal} \\
    -\hat{\bm{l}} & \text{if } \bm{e}_B \text{ is distal}
  \end{cases}
\end{equation}

Then the sign is checked so that $\hat{\bm{z}}$ points \emph{away from} the
distal end:

\begin{equation}
  \hat{\bm{z}} =
  \begin{cases}
    -\hat{\bm{z}}_{\mathrm{raw}}
      & \text{if } \hat{\bm{z}}_{\mathrm{raw}} \cdot
        (\bm{e}_{\mathrm{distal}} - \bar{\bm{p}}) > 0 \\
    \hat{\bm{z}}_{\mathrm{raw}} & \text{otherwise}
  \end{cases}
\end{equation}

\subsection{Anterior direction via surface-normal sector analysis}

An orthonormal reference frame perpendicular to $\hat{\bm{z}}$ is built using
Gram-Schmidt:

\begin{equation}
  \hat{\bm{b}}_x = \mathrm{normalise}\!\left(
    \bm{t} - (\bm{t} \cdot \hat{\bm{z}})\,\hat{\bm{z}}
  \right),
  \qquad \bm{t} =
  \begin{cases}
    (1,0,0)^\top & \text{if } |\hat{z}_x| < 0.9 \\
    (0,1,0)^\top & \text{otherwise}
  \end{cases}
\end{equation}

\begin{equation}
  \hat{\bm{b}}_y = \hat{\bm{z}} \times \hat{\bm{b}}_x
\end{equation}

For each vertex normal $\hat{\bm{n}}_i$, the component perpendicular to
$\hat{\bm{z}}$ is:

\begin{equation}
  \bm{n}_i^\perp = \hat{\bm{n}}_i - (\hat{\bm{n}}_i \cdot \hat{\bm{z}})\,\hat{\bm{z}}
\end{equation}

This is projected into the 36-sector angular histogram:

\begin{equation}
  \theta_i = \mathrm{atan2}\!\left(
    \bm{n}_i^\perp \cdot \hat{\bm{b}}_y,\;
    \bm{n}_i^\perp \cdot \hat{\bm{b}}_x
  \right), \qquad \theta_i \in (-\pi, \pi]
\end{equation}

\begin{equation}
  s_i = \left\lfloor \frac{\theta_i + \pi}{2\pi} \cdot 36 \right\rfloor
        \bmod 36
\end{equation}

Each sector $s$ accumulates the mean perpendicular normal magnitude:

\begin{equation}
  \mathrm{score}(s) = \frac{1}{|s|} \sum_{i\,:\,s_i = s} \|\bm{n}_i^\perp\|
\end{equation}

The sector with the highest score gives the raw anterior angle:

\begin{equation}
  s^* = \arg\max_s\; \mathrm{score}(s),
  \qquad
  \theta^* = \frac{s^* + 0.5}{36} \cdot 2\pi - \pi
\end{equation}

\begin{equation}
  \hat{\bm{x}}_{\mathrm{raw}} =
    \cos\theta^*\;\hat{\bm{b}}_x + \sin\theta^*\;\hat{\bm{b}}_y
\end{equation}

\subsection{Final orthonormal basis}

The raw anterior direction is orthogonalised against $\hat{\bm{z}}$ (to remove
any residual component along the long axis), then the medial axis is derived by
the right-hand rule:

\begin{equation}
  \hat{\bm{x}} = \mathrm{normalise}\!\left(
    \hat{\bm{x}}_{\mathrm{raw}}
    - (\hat{\bm{x}}_{\mathrm{raw}} \cdot \hat{\bm{z}})\,\hat{\bm{z}}
  \right)
\end{equation}

\begin{equation}
  \hat{\bm{y}} = \hat{\bm{z}} \times \hat{\bm{x}}
\end{equation}

The resulting frame uses the convention:

\begin{center}
\begin{tabular}{ll}
  \toprule
  Axis & Direction \\
  \midrule
  $+X$ & Anterior \\
  $+Y$ & Medial   \\
  $+Z$ & Proximal (up) \\
  \bottomrule
\end{tabular}
\end{center}

% ─────────────────────────────────────────────────────────────────────────────
\section{Orientation Transform}

\subsection{Rotation matrix}

The three axis vectors are assembled as columns of a rotation matrix:

\begin{equation}
  \bm{R} = \bigl[\,\hat{\bm{x}} \;\big|\; \hat{\bm{y}} \;\big|\; \hat{\bm{z}}\,\bigr]
  \in SO(3)
\end{equation}

Because $\bm{R}$ is orthonormal, its inverse is its transpose:
$\bm{R}^{-1} = \bm{R}^\top$.

\subsection{Full 4$\times$4 homogeneous transform}

The final transform translates the distal end to the origin and then applies
the inverse rotation to align the limb with the world axes:

\begin{equation}
  \bm{T}_{-\bm{e}} =
  \begin{pmatrix}
    1 & 0 & 0 & -e_x \\
    0 & 1 & 0 & -e_y \\
    0 & 0 & 1 & -e_z \\
    0 & 0 & 0 & 1
  \end{pmatrix},
  \qquad \bm{e} = \bm{e}_{\mathrm{distal}}
\end{equation}

\begin{equation}
  \boxed{\bm{M} = \bm{R}^\top \bm{T}_{-\bm{e}}}
\end{equation}

\subsection{Applying the transform to vertices}

Each vertex position is transformed in homogeneous coordinates:

\begin{equation}
  \bm{p}_i' =
  \bm{M}
  \begin{pmatrix} \bm{p}_i \\ 1 \end{pmatrix}
  = \bm{R}^\top (\bm{p}_i - \bm{e}_{\mathrm{distal}})
\end{equation}

After transformation: the distal tip sits at the origin $(0,0,0)$, the
proximal end is along $+Z$, and the anterior face points toward $+X$.

\subsection{Anterior correction (manual)}

When the user clicks \textit{Anterior Facing Me}, the camera's current angle
$\phi$ in the $XY$-plane is used to apply a further rotation about $Z$:

\begin{equation}
  \bm{R}_Z(\phi) =
  \begin{pmatrix}
    \cos\phi & -\sin\phi & 0 \\
    \sin\phi &  \cos\phi & 0 \\
    0        & 0         & 1
  \end{pmatrix}
\end{equation}

% ─────────────────────────────────────────────────────────────────────────────
\section{Mesh Volume}

The volume of the largest connected body is computed via the
\emph{signed-tetrahedron} (divergence theorem) method.  For each triangular
face with vertices $\bm{a}, \bm{b}, \bm{c}$, the signed scalar triple product
gives the volume of the tetrahedron formed with the origin:

\begin{equation}
  V_f = \bm{a} \cdot (\bm{b} \times \bm{c})
      = a_x(b_y c_z - b_z c_y)
      + b_x(c_y a_z - c_z a_y)
      + c_x(a_y b_z - a_z b_y)
\end{equation}

Summing over all faces belonging to the largest connected component and taking
the absolute value:

\begin{equation}
  V = \frac{1}{6}\left|\sum_{f} V_f\right|
  \quad \text{(mm}^3\text{)}
\end{equation}

% ─────────────────────────────────────────────────────────────────────────────
\section{Marker Detection}

\subsection{Shape signature}

The same PCA and OBB (Oriented Bounding Box) procedure described in
Section~1 is applied to the marker geometry to extract three sorted extents
$d_1 \geq d_2 \geq d_3$ along the PCA axes.  Two dimensionless ratios
characterise the shape:

\begin{equation}
  \kappa = \frac{d_3}{d_1}
  \quad\text{(compactness: 0 = rod-like, 1 = cubic)}
  \qquad
  \rho = \frac{d_2}{d_1}
  \quad\text{(mid-axis ratio)}
\end{equation}

\subsection{Spatial grid search}

The scan is divided into a regular grid with cell size $S = d_1$ (the
marker's largest dimension).  For each occupied grid cell, a
$(2g+1)^3$-block neighbourhood (default $g = 1$, i.e.\ $3^3 = 27$ cells)
is gathered and scored:

\begin{equation}
  \text{score} =
    0.45\,\bigl(1 - |\kappa_c - \kappa_m|\bigr)
  + 0.25\,\bigl(1 - |\rho_c   - \rho_m|\bigr)
  + 0.30\,\exp\!\left(-\frac{|d_{1,c} - d_{1,m}|}{d_{1,m}}\right)
\end{equation}

where subscript $c$ denotes the candidate neighbourhood and $m$ the marker
reference signature.  The neighbourhood with the highest score (above a
minimum threshold of 0.25) is declared the detected marker location.

\subsection{Highlight radius}

The highlighted vertex set shown in the 3-D viewer includes all scan vertices
within a sphere of radius $r_h = 0.65\,d_{1,m}$ centred on the detected
centroid.

% ─────────────────────────────────────────────────────────────────────────────
\section*{Summary of key equations}

\begin{align}
  \bar{\bm{p}}    &= \tfrac{1}{n}\textstyle\sum_i \bm{p}_i
  \tag{Centroid} \\[4pt]
  \bm{C}          &= \tfrac{1}{n}\textstyle\sum_i
    (\bm{p}_i-\bar{\bm{p}})(\bm{p}_i-\bar{\bm{p}})^\top
  \tag{Covariance} \\[4pt]
  \hat{\bm{l}}    &= \lim_{k\to\infty}
    \tfrac{\bm{C}^k\bm{v}_0}{\|\bm{C}^k\bm{v}_0\|}
  \tag{Long axis} \\[4pt]
  \hat{\bm{x}},\,\hat{\bm{y}},\,\hat{\bm{z}}
                  &= \text{Gram-Schmidt from } \hat{\bm{l}}
  \tag{Basis} \\[4pt]
  \bm{M}          &= \bm{R}^\top \bm{T}_{-\bm{e}_{\mathrm{distal}}}
  \tag{Orient transform} \\[4pt]
  V               &= \tfrac{1}{6}\bigl|\textstyle\sum_f
    \bm{a}_f\cdot(\bm{b}_f\times\bm{c}_f)\bigr|
  \tag{Volume}
\end{align}

\end{document}
